\documentclass[../../../include/open-logic-section]{subfiles}

\begin{document}

\olfileid{sth}{cardinals}{hp}
\olsection{Lampiran: Prinsip Hume}

Ing \olref[cp]{sec}, kita wis njlentrehaké Prinsip Cantor.
Prinsip mau yaiku:
\begin{align*}
	\card{A} = \card{B} & \text{ yen lan mung yen } A \approx B.
\intertext{Iki mirip banget karo kang saiki diarani
\emph{Prinsip Hume}, kang kandha:}
	\fregenum{x} {F(x)} = \fregenum{x}{G(x)} & \text{ yen lan mung yen } F \sim G
\end{align*}
ing kéné “$F \sim G$” nyingkat pratelan yèn cacahé $F$ persis
padha karo cacahé $G$, yaiku objek-objek $F$ bisa dipasangake
kanthi bijeksi karo objek-objek $G$, yaiku:
\begin{align*}
	\exists R(&\forall v\forall y(Rvy \lif (Fv \land Gy)) \land {}\\
		&\forall v(Fv \lif \lexists![y][Rvy]) \land {}\\
		&\forall y(Gy \lif \lexists![v][Rvy]))
\end{align*}
Nanging Prinsip Hume lan Prinsip Cantor béda jinisé. Ing
rumusan Prinsip Cantor, variabel “$A$” lan “$B$” iku
term orde kapisan kang nuduhaké \emph{himpunan}. Ing
rumusan Prinsip Hume, “$F$”, “$G$”, lan “$R$”
\emph{dudu} term orde kapisan; kabèh mau dumunung ing
\emph{panggonan predikat}. (Mbokmenawa nuduhaké
\emph{sipat}.) Mula Prinsip Hume bisa kita jlentrehaké
mangkéné: cacahé $F$ padha karo cacahé $G$ yen lan mung yen
objek-objek $F$ bisa dibijeksèkaké karo~$G$. Prinsip iki
diarani \emph{Prinsip Hume}, amarga Hume tau nulis:
\begin{quote}
  Nalika rong wilangan digandhèngaké saéngga saben satuan ing
  siji tansah nduwèni pasangan satuan ing sijiné, kita nyebut
  loro-loroné padha.
  \citep[Pt.III Bk.1 \S1]{Hume1740}
\end{quote}
Prinsip Hume banjur diangkat dadi wigati ing logika
matematika kontemporer déning \citet[\S63]{Frege1884}.
Dheweke ngutip pethikan Hume mau sadurungé (ing intiné)
nggambarake apa kang kita sebut Prinsip Hume.

Wigatèkna kamiripan struktur antara Prinsip Hume lan
Hukum Dhasar~V. Kita ngrumusaké ing \olref[story][blv]{sec}
mangkéné:
\[
	\fregeext{x}{F(x)} = \fregeext{x}{G(x)} \text{yen lan mung yen } \lforall[x][(F(x) \liff G(x))].
\]
Ing titik iki, sawetara panjlentrehan lan pambandhing bisa migunani.

Ana rong cara mangertèni prinsip kaya Prinsip Hume utawa Hukum
Dhasar~V: kanthi \emph{predikatif} utawa
\emph{impredikatif} (élinga
\olref[story][predicative]{sec}). Miturut pamacaan impredikatif
marang Hukum Dhasar~V, kanggo saben~$F$, objek
$\fregeext{x}{F(x)}$ kalebu ing domain kuantifikasi kang
dienggo kanggo ngrumusaké Hukum Dhasar~V iku dhéwé.
Semono uga, miturut pamacaan impredikatif marang Prinsip
Hume, kanggo saben~$F$, objek $\fregenum{x}{F(x)}$
kalebu ing domain kuantifikasi kang dienggo kanggo ngrumusaké
Prinsip Hume. Kosokbaliné, miturut pangertèn
\emph{predikatif}, objek $\fregeext{x}{F(x)}$ lan~$\fregenum{x}{F(x)}$
iku anggota domain kang \emph{béda}.

Yèn Hukum Dhasar~V diwaca kanthi impredikatif, bakal
nuwuhaké inkonsistensi liwat Komprehensi Naif (rinciné
delengen \olref[story][blv]{sec}). Kaya Komprehensi Naif,
prinsip iki bisa digawe konsisten kanthi pamacaan
\emph{predikatif}. Nanging bisa uga ora nyukupi kabèh
kang kita karepake.

Prinsip Hume, kosokbaliné, \emph{bisa} diwaca kanthi
impredikatif tanpa inkonsistensi. Kanthi pamacaan mangkono,
prinsip iki cukup kuwat.

Kanggo nggambarake: gatekna predikat “$x \neq x$”,
kang cetha ora disukupi déning objek apa waé. Prinsip
Hume banjur menehi sawijining objek $\# x( x\neq
x)$. Kita bisa nganggep iki minangka wilangan~$0$. Saiki,
miturut pangertèn \emph{impredikatif}---nanging
\emph{mung} miturut pangertèn impredikatif---objek $0$ iki
kalebu ing domain kuantifikasi wiwitan kita. Mula Prinsip
Hume bisa ditrapaké marang predikat “$x = 0$” kanggo
ngasilaké objek $\#x (x =
0)$. Kita bisa nganggep iki minangka wilangan~$1$.
Kajaba iku, Prinsip Hume ngakibataké $0 \neq 1$,
amarga ora mungkin ana bijeksi saka objek kang ora padha
karo dhiriné menyang objek kang padha karo $0$
(sing kapisan ora ana, déné sing kapindho ana siji).
Saiki, miturut pamacaan impredikatif manèh, $1$~kalebu
ing domain kuantifikasi wiwitan kita. Mula Prinsip Hume
bisa ditrapaké marang predikat “$(x = 0 \lor x = 1)$”
kanggo ngasilaké objek $\#x(x = 0 \lor
x = 1)$. Kita bisa nganggep iki minangka wilangan~$2$,
lan bisa mbuktèkaké $0\neq 2$ lan $1 \neq 2$, lan saterusé.

Dadi, yèn diwaca kanthi impredikatif, Prinsip Hume
ngakibataké anané \emph{objek tanpa wates cacahé}.
Iki nyurung para \emph{logisis neo-Fregean} kanggo
nganggo Prinsip Hume minangka dhasar aritmetika.

Nanging Frege \emph{dhéwé} ora nganggo Prinsip Hume
minangka dhasar aritmetika. Frege malah mbuktèkaké
Prinsip Hume saka sawijining definisi eksplisit:
$\fregenum{x}{F(x)}$ ditetepaké minangka ekstensi
konsep $F \sim \Phi$. Ing istilah modern, kita bisa
nyoba nulis iki minangka
$\fregenum{x}{F(x)} = \Setabs{G}{F \sim G}$;
nanging iki mbalèkaké kita marang masalah Komprehensi Naif.

\end{document}
